\section{Ingeniería de software de investigación}

Software es un programa de computo y documentación asociada, la ingeniería de software es una disciplina de la ingeniería que se interesa en todos los aspectos de la producción del software en todas sus etapas. Este proceso incluye cuatro actividades fundamentales las cuales son:

\begin{itemize}
	\item \textbf{Especificación del software}: Incluye la toma de requerimientos, definición el alcance y limitaciones y especificación de casos de uso.  
	
	\item \textbf{Desarrollo de software}: Diseño y programación del software.
	
	\item \textbf{Validación de software}: Se verifica que el software cumpla con los requerimientos establecidos.
	
	\item \textbf{Evolución del software}: Cambios en el software que refleje los cambios o añadidos en los requerimientos.
\end{itemize}

Existen múltiples tipos de aplicaciones de ingeniería de software, como aplicaciones independientes, sistemas embebidos, de entretenimiento y muy particularmente para este caso, sistemas de modelado y simulación. Estos son utilizados para modelar situaciones físicas u objetos que interactúen de forma dinámica. Estos sistemas suelen ser computacionalmente intensivos y requiere un alto desempeño de procesos en paralelo \cite{ingenieria_1}.
